% EXCERPT of VCFv4.5.tex from https://github.com/samtools/hts-specs (master @ 510c107ba89aa13113818c9e4fa3071283cb7f7b), retrieved 2026-11-17.
% Kept: lines 40-71 (example file), 103-148 (meta-information lines), 291-330 (##assembly and ##contig), 2147-2166 (BCF contig dictionary). Nothing else edited.
% ---- lines 40-71 ----
\subsection{An example}
\scriptsize
\begin{verbatim}
##fileformat=VCFv4.5
##fileDate=20090805
##source=myImputationProgramV3.1
##reference=file:///seq/references/1000GenomesPilot-NCBI36.fasta
##contig=<ID=20,length=62435964,assembly=B36,md5=f126cdf8a6e0c7f379d618ff66beb2da,species="Homo sapiens",taxonomy=x>
##phasing=partial
##INFO=<ID=NS,Number=1,Type=Integer,Description="Number of Samples With Data">
##INFO=<ID=DP,Number=1,Type=Integer,Description="Total Depth">
##INFO=<ID=AF,Number=A,Type=Float,Description="Allele Frequency">
##INFO=<ID=AA,Number=1,Type=String,Description="Ancestral Allele">
##INFO=<ID=DB,Number=0,Type=Flag,Description="dbSNP membership, build 129">
##INFO=<ID=H2,Number=0,Type=Flag,Description="HapMap2 membership">
##FILTER=<ID=q10,Description="Quality below 10">
##FILTER=<ID=s50,Description="Less than 50% of samples have data">
##FORMAT=<ID=GT,Number=1,Type=String,Description="Genotype">
##FORMAT=<ID=GQ,Number=1,Type=Integer,Description="Genotype Quality">
##FORMAT=<ID=DP,Number=1,Type=Integer,Description="Read Depth">
##FORMAT=<ID=HQ,Number=2,Type=Integer,Description="Haplotype Quality">
#CHROM POS     ID        REF    ALT     QUAL FILTER INFO                              FORMAT      NA00001        NA00002        NA00003
20     14370   rs6054257 G      A       29   PASS   NS=3;DP=14;AF=0.5;DB;H2           GT:GQ:DP:HQ 0|0:48:1:51,51 1|0:48:8:51,51 1/1:43:5:.,.
20     17330   .         T      A       3    q10    NS=3;DP=11;AF=0.017               GT:GQ:DP:HQ 0|0:49:3:58,50 0|1:3:5:65,3   0/0:41:3
20     1110696 rs6040355 A      G,T     67   PASS   NS=2;DP=10;AF=0.333,0.667;AA=T;DB GT:GQ:DP:HQ 1|2:21:6:23,27 2|1:2:0:18,2   2/2:35:4
20     1230237 .         T      .       47   PASS   NS=3;DP=13;AA=T                   GT:GQ:DP:HQ 0|0:54:7:56,60 0|0:48:4:51,51 0/0:61:2
20     1234567 microsat1 GTC    G,GTCT  50   PASS   NS=3;DP=9;AA=G                    GT:GQ:DP    0/1:35:4       0/2:17:2       1/1:40:3
\end{verbatim}
\normalsize
This example shows (in order): a good simple SNP, a possible SNP that has been filtered out because its quality is below 10, a site at which two alternate alleles are called, with one of them (T) being ancestral (possibly a reference sequencing error), a site that is called monomorphic reference (i.e.\ with no alternate alleles), and a microsatellite with two alternative alleles, one a deletion of 2 bases (TC), and the other an insertion of one base (T).
Genotype data are given for three samples, two of which are phased and the third unphased, with per sample genotype quality, depth and haplotype qualities (the latter only for the phased samples) given as well as the genotypes.
The microsatellite calls are unphased.
% ---- lines 103-148 ----
\subsection{Meta-information lines}
File meta-information lines start with ``\verb|##|'' and must appear first in the VCF file, before the header line (section~\ref{header-line}) and data record lines (section~\ref{data-lines}).
They may be either \emph{unstructured} or \emph{structured}.

An \emph{unstructured} meta-information line consists of a~\emph{key} (denoting the type of meta-information recorded) and a~\emph{value} (which may not be empty and must not start with a `\verb|<|' character), separated by an `\verb|=|' character:
\begin{quote}
\verb|##|\emph{key}\verb|=|\emph{value}
\end{quote}
Several unstructured meta-information lines are defined in this specification, notably \verb|##fileformat|.
Others not defined by this specification, e.g.\ \verb|##fileDate| and \verb|##source|, are commonly found in VCF files.
These typically have meanings that are obvious, or they are immaterial for processing the file, or both.

A \emph{structured} meta-information line is similar, but the value is itself a comma-separated list of key=value pairs, enclosed within `\verb|<|' and `\verb|>|' characters:
\begin{quote}
\verb|##|\emph{key}\verb|=<|\emph{key}\verb|=|\emph{value}\verb|,|\emph{key}\verb|=|\emph{value}\verb|,|\emph{key}\verb|=|\emph{value}\verb|,|\ldots\verb|>|
\end{quote}
All structured lines require an ID which must be unique within their type, i.e., within all the meta-information lines with the same ``\verb|##|\emph{key}\verb|=|'' prefix.
For all of the structured lines (\verb|##INFO|, \verb|##FORMAT|, \verb|##FILTER|, etc.) described in this specification, optional fields can be included.
For example:
\begin{verbatim}
##INFO=<ID=ALLELEID,Number=A,Type=String,Description="Allele ID",Source="ClinVar",Version="20220804">
\end{verbatim}
In the above example, the optional fields of ``Source'' and ``Version'' are provided.
The values of optional fields must be written as quoted strings, even for numeric values.
Other structured lines not defined by this specification may also be used; the only required field for such lines is the required \verb|ID| field.

It is recommended in VCF and required in BCF that the header includes tags describing the reference and contigs backing the data contained in the file.
These tags are based on the SQ field from the SAM spec; all tags are optional (see the VCF example above).

To aid human readability, the order of fields should be ID, Number, Type, Description, then any optional fields.
Implementations must not rely on the order of the fields within structured lines and are not required to preserve field ordering.

Meta-information lines are optional, but if they are present then they must be completely well-formed.
Other than \verb|##fileformat|, they may appear in any order.
Note that BCF, the binary counterpart of VCF, requires that all entries are present.
It is recommended to include meta-information lines describing the entries used in the body of the VCF file.


\subsubsection{File format}
A single `fileformat' line is always required, must be the first line in the file, and details the VCF format version number.
For VCF version 4.5, this line is:

\begin{verbatim}
##fileformat=VCFv4.5
\end{verbatim}

% ---- lines 291-330 ----
\subsubsection{Assembly field format}
Breakpoint assemblies for structural variations may use an external file:
\begin{verbatim}
##assembly=url
\end{verbatim}

The URL field specifies the location of a fasta file containing breakpoint assemblies referenced in the VCF records for structural variants via the BKPTID INFO key.

\subsubsection{Contig field format}
\label{sec-contig-field}
It is recommended for VCF, and required for BCF, that the header includes tags describing the contigs referred to in the file.
The structured \texttt{contig} field must include the ID attribute and can include additional optional attributes with
the following ones reserved:
\begin{itemize}
  \item length: the length of the sequence
  \item md5: MD5 checksum of the sequence as defined in the Sam specification v1\footnote{See Reference MD5 calculation
      section in \href{https://samtools.github.io/hts-specs/SAMv1.pdf}{\tt SAM Format Specification}.} Briefly, the digest
      is calculated excluding all characters outside of the inclusive range 33 (`\char33') to 126 (`\char126').
      and all lowercase characters converted to uppercase. The MD5 digest is calculated as described in
      \href{https://tools.ietf.org/html/rfc1321}{\sl RFC 1321} and presented as a 32 character lowercase hexadecimal number.
  \item URL: tag to indicate where the sequence can be found
\end{itemize}

For example:
{\scriptsize
\begin{verbatim}
##contig=<ID=ctg1,length=81195210,URL=ftp://somewhere.example/assembly.fa,md5=f126cdf8a6e0c7f379d618ff66beb2da,...>
\end{verbatim}
}
\noindent
Contig names follow the same rules as the SAM format's reference sequence names:
they may contain any printable ASCII characters in the range \verb|[!-~]| apart from `{\tt\verb|\|\,,\,"`'\,()\,[]\,\verb|{}|\,<>}' and may not start with `{\tt *}' or `{\tt =}'.
Thus they match the following regular expression:
\begin{verbatim}
    [0-9A-Za-z!#$%&+./:;?@^_|~-][0-9A-Za-z!#$%&*+./:;=?@^_|~-]*
\end{verbatim}
\noindent
In particular, excluding commas facilitates parsing \verb|##contig| lines, and excluding the characters `\verb|<>[]|' and initial~`{\tt *}' avoids clashes with symbolic alleles.
The contig names must not use a reserved symbolic allele name.

% ---- lines 2147-2166 ----
\subsubsection{Dictionary of contigs}

The CHROM field in BCF2 is encoded as an integer offset into the list of \verb|##contig| field headers in the VCF header.
The offsets begin, like the dictionary of strings, at 0.
So for example if in BCF2 the contig value is 10, this indicates that the actual chromosome is the 11th element in the ordered list of \verb|##contig| elements.
Here's a more concrete example:

\small
\begin{verbatim}
##contig=<ID=20,length=62435964,assembly=B36,md5=f126cdf8a6e0c7f379d618ff66beb2da,species="Homo sapiens">
##contig=<ID=21,length=46944323,assembly=B36,md5=f1b74b7f9f4cdbaeb6832ee86cb426c6,species="Homo sapiens">
##contig=<ID=22,length=49691432,assembly=B36,md5=2041e6a0c914b48dd537922cca63acb8,species="Homo sapiens">
#CHROM POS ID REF ALT QUAL FILTER INFO
20 1 . T A . PASS .
21 2 . T A . PASS .
22 3 . T A . PASS .
\end{verbatim}
\normalsize

the actual CHROM field values in the encoded BCF2 records would be 0, 1, and 2 corresponding to the first (offset 0) \verb|##contig| element, etc.
